% Pursuing Stacks (cotes 134-2 to 134-8): archive cote 134-6, site page 109 (PDF page 110), typescript page 427.
% typescript. Draft compared: ps5.tex:7173-7240.
% AI-transcribed by Claude Opus 5.5 (claude-opus-5-5, adaptive thinking, effort high), Message Batches API, batch ps3 (msgbatch_01XdFmuVVBAcgc7EFxdtCqsj), 2026-11-12; not yet checked by eye.
\scanpage{134-6/109}
\mspage{427}

viewed as being just the invariants $\mathrm{Ext}^{i}(\underline{Z}_{\hand{\mathcal{X}}}\hand{,}F)$ \unsure{in} the category of all
abelian sheaves, where $\underline{Z}_{\hand{\mathcal{X}}}$ is the constant sheaf on $\hand{\mathcal{X}}$ with value $\underline{Z}$ .
The similar \xout{statement} \inserted{fact} holds when F is any sheaf of modules over a sheaf
of rings $\underline{O}_{\hand{\mathcal{X}}}$ on $\hand{\mathcal{X}}$ , with $\underline{Z}_{\hand{\mathcal{X}}}$ being replaced by $\underline{O}_{\hand{\mathcal{X}}}$ in the interpretation
above :
\[ \mathrm{H}^{i}(\hand{\mathcal{X}}, F) \hand{\simeq} \mathrm{Ext}^{i}_{\underline{O}_{\hand{\mathcal{X}}}}(\underline{O}_{\hand{\mathcal{X}}}, F) \quad , \]
which is often quite useful in the cohomology formalism. We are going
to restrict here to the case of a constant sheaf of rings, defined by a
\inserted[typed]{fixed} commutative ring k, which will allow us to play around as announced
with the duality relation between A and B=A$^{\mathrm{o}}$ , provided moreover we
take more general coefficients still, namely \deleted{coefficients} \inserted[typed]{(pre)sheaves with values} in a\deleted{ny} \hand{given}
k-additive category \inserted{$\mathcal{M}$} stable under the relevant limits. (Presumably the
case of a locally constant commutative sheaf of rings could be dealt
with too, but we'll not dive into this here ! ). Another important
(and by now familiar ??) conceptual point is that, \deleted{more} \xout{important} \inserted{rather} \inserted[typed]{relevant} than \hand{\unsure{in}}
the Ext$^{i}$ 's which give only partial information, we are interested in
the object \xout{of a suitable derived category} they come from(as the ``cohomo-%
logy objects''), namely \inserted{the} object\unsure{\deleted{s}|s} $\mathrm{RHom}_{\hand{\underline{O}_{\mathcal{X}}}}(\hand{L,}F)$ in a suitable derived category\hand{\unsure{.}}
In the present case when $\hand{\mathcal{X}}$ is the topos associated to the small category
A, hence the category of $\underline{O}_{\hand{\mathcal{X}}}$\xout{modules}-Modules has sufficiently many projectives
(namely direct sums of sheaves of the type $\underline{O}_{\hand{\mathcal{X}}}^{(a)}$ with a in A), the
$\mathrm{RHom}_{\underline{O}_{\hand{\mathcal{X}}}}(L,F)$ may be computed, taking a projective resolution L. of E ,
by the formula
\[ \mathrm{RHom}_{\underline{O}}(L,F) \hand{\simeq} \mathrm{Hom}^{\bullet}_{\underline{O}}(L. \ , F) \]
( an isomorphism in D$^{+}$(\deleted{\illegible{1}}k-Mod) \inserted{say} ) , and similarly when replacing L and
F by arguments L. and F$^{\bullet}$ in D$^{-}$ and D$^{+}$ of the category of $\underline{O}_{\hand{\mathcal{X}}}$-Modules.
As a result, we get a pairing, computable \inserted{here} using projective resolution\hand{s}
of the argument L. :
\[ \hand{(*)} \qquad (L. \ , F^{\bullet}) \hand{\mapsto} \mathrm{RHom}_{\underline{O}_{\hand{\mathcal{X}}}}(L.,F^{\bullet}) \ : \ D^{-}(\underline{O}_{\hand{\mathcal{X}}}) \times D^{+}(\underline{O}_{\hand{\mathcal{X}}}) \hand{\longrightarrow} D^{+}(\text{k-Mod}) \ , \]
where k is a commutative ring and $\underline{O}_{\hand{\mathcal{X}}}$ is endowed with a structure of
k-algebra. Using $\underline{\mathrm{Hom}}_{\hand{\underline{O}_{\mathcal{X}}}}$ and it's total derived functor, we get likewise
\[ \hand{(**)} \qquad (L.,F^{\bullet}) \hand{\mapsto} \mathrm{R}\underline{\mathrm{Hom}}_{\underline{O}_{\hand{\mathcal{X}}}}(L.,F^{\bullet}) \ : \ D^{-}(\underline{O}_{\hand{\mathcal{X}}}) \times D^{+}(\underline{O}_{\hand{\mathcal{X}}}) \hand{\longrightarrow} D^{+}(\underline{O}_{\hand{\mathcal{X}}}) \ , \]
with
\[ \underline{\mathrm{RHom}}_{\hand{\underline{O}_{\mathcal{X}}}}(L.,F^{\bullet}) \hand{\simeq} \underline{\mathrm{Hom}}_{\underline{O}_{\hand{\mathcal{X}}}}^{\bullet\bullet}(\hand{\mathcal{L}}. \ , F^{\bullet}) \quad , \]
where $\hand{\mathcal{L}}$. is a projective resolution of L. , and $\underline{\mathrm{Hom}}^{\bullet\bullet}$ stands for
the simple complex associated to the double complex obtained by taking

\margin{109}
